\centerline{\bf \Huge ДЗ. Процессы и операции}
\centerline{\bf \Huge }

\begin{flushright}
        \scriptsize{\textbf{Разрешаю не доказывать Оценку во 2-ом и в 4-ом задании. Но если сможете, то получите дополнительные баллы!}}
\end{flushright}

\problem{\textbf{1}} В очереди в столовую ЭТЛ стоят 2025 школьников так, что
мальчики и девочки среди них чередуются. (Первым стоит мальчик, за ним — девочка, за
ней — снова мальчик и так далее.) Любой мальчик, за спиной которого в очереди стоит девочка, может
поменяться с ней местами. Через некоторое время оказалось, что все девочки стоят в начале
очереди, а все мальчики — в конце. Сколько обменов было совершено?

\problem{\textbf{2}} Начальник I умеет умножать числа на 3, Начальник II — прибавлять 1, Начальник III — умножать число само на себя (возводить во вторую степень), Начальник IV — вычитать 2. В каком порядке им нужно выполнять свои операции (не обязательно каждую и не обязательно по одному разу), чтобы как можно быстрее получить из числа 1 число 2025?

\problem{\textbf{3}} Из последовательности натуральных
чисел 1, 2, 3, . . . удалили все точные квадраты (целые числа возведённые во вторую степень 1, 4, 9, 16, 25, ...) и все точные кубы (целые числа возведённые в третью степень 1, 8, 27, 64, 125, ...). Какое число будет
находиться на 2025 месте среди оставшихся?

\problem{\textbf{4}} У Начальника есть 5 пачек с юанями по 100 штук в каждой. К нему
подошли пятеро учеников. Первому ученику нужно 120 купюр, второму и третьему — по
110, четвёртому 80, а пятому 50. Как Начальнику отсчитать каждому ученику нужное число купюр за наименьшее количество времени, если за одну секунду он отсчитывает ровно одну купюру? Сколько секунд он на это потратит?




